\section*{Part VI: The Industry Atlas: Different Mechanisms for Different Markets}

\subsection*{64. FMCG: Recognition, Access, and Repeat Behavior}

A fast-moving consumer brand competes in a world of scarce shelf space, fast decisions, retailer incentives, and frequent repurchase. A customer might never articulate a complex reason for choosing 1 biscuit over another. Availability, learned taste, brand cues, pack sizes, price points, retailer recommendations, and habit dominate.

A small FMCG brand might differentiate by serving a precise eating occasion, distinct flavor preference, credible ingredient proposition, easier refill behavior, or superior access to a niche channel. But product-market fit cannot be inferred from initial trials fueled by promotions. Measure rate of sale per outlet, repeat rate, contribution after retailer margins and trade spend, availability, returns, and product experience. In this industry, being recognized and conveniently bought is often as important as being technically different.

\subsection*{65. Luxury and Symbolic Goods: Purchasing Recognized Meaning}

Luxury categories combine craftsmanship, design, rarity, provenance, cultural recognition, and social interpretation. The value of an object is partly positional: it matters that a particular community recognizes its history or scarcity. Quality is important without explaining the entire price.

A newcomer cannot fabricate decades of collective meaning. It can begin with remarkable craft, coherent design, authentic scarcity from capacity constraints, truthful provenance, and a community that cares about the work. High prices and elegant photography alone do not create a durable luxury brand. The key test is whether customers willingly pay for the whole experience and continue to value the object after the novelty fades.

\subsection*{66. Software-as-a-Service: Workflow Depth and Adoption Economics}

B2B software features tend to converge. A feature that was novel last quarter may soon be available as an extension of an incumbent platform. The better question is whether the tool owns or substantially improves a consequential workflow.

For example, ``accurate meeting summaries'' is easy to compare and imitate. ``Approved client meeting records with review status, traceability, role controls, and integration into the case workflow'' is a different unit of value. The latter imposes additional development, security, customer-support, and implementation burdens. A position must be tested with the real buying committee, not merely admired by an individual user.

\subsection*{67. AI Products: Do Not Sell Commodity Model Access as Uniqueness}

When a large share of a product is powered by an external or widely available foundation model, its capabilities can improve or commoditize rapidly. The mere presence of ``AI,'' chat, retrieval, or an agent framework is rarely a secure position. A defensible differentiator lies in domain workflow, proprietary lawfully accessible context, evaluation infrastructure, user trust, integration, technical reliability, and mechanisms that reduce expensive mistakes.

A founder must explicitly ask: \textbf{If a frontier model becomes twice as capable and half as costly, does our product become more valuable, less valuable, or redundant?} If customers can reproduce the same result using a generic assistant and a prompt, your value is convenience and distribution rather than unique intelligence.

\subsection*{68. Consumer Apps: Relevance and Habit Compete with Novelty}

A consumer application must win a place in attention and routine. The first adoption may result from novelty, but durable usage depends on repeated utility, entertainment, social interaction, or identity. A new app can be functionally sophisticated and still fail because nobody remembers it or wants to change existing habits.

Differentiate by a compelling use occasion, rapid time to value, a highly relevant community, or a neglected behavior that current tools support badly. Measure retention by cohort and actual meaningful use, not registrations.

\subsection*{69. Marketplaces: The Product Is Often Credible Liquidity}

A marketplace may appear to sell matching software while its real product is \textbf{the probability of a suitable transaction at acceptable time, cost, and risk.} Two-sided growth is not automatically a network effect: participants must be relevant to 1 another, supply must be reliable, and transactions must complete.

A strong beachhead can be a narrow geographic cell, specialty category, recurring demand occasion, or constrained set of trusted suppliers. Differentiation comes from verification, service levels, matching quality, dispute resolution, integration, or density.

\subsection*{70. Professional Services: Judgment and Accountability}

Consulting firms often claim they are strategic, innovative, customer-focused, and expert. These generic descriptions are nearly impossible for buyers to compare. A more meaningful position combines a specific problem, client context, measurable deliverable, scope boundaries, competent personnel, and reliable execution.

Consider an insurer who needs an auditable evaluation of a claims-triage model. A narrowly defined assessment with transparent criteria and documented limitations is easier to buy and more useful than a generic ``AI transformation'' project. Measure profitability per \textbf{constrained expert hour}, rework, project delays, scope creep, referrals, and repeat demand.

\subsection*{71. Manufacturing and Supply: Reliability Can Dominate Features}

A manufacturer competes not only on product specifications but yield, tolerances, documentation, supply continuity, defect rate, lead time, maintenance, and compatibility with customers' production processes. Buyers pay more to avoid an expensive production stoppage even when the item itself appears commodity-like.

Differentiation arises from application engineering, shorter replenishment cycles, traceability, reliable stock, and predictable technical support. But physical inventory and quality commitments impose working-capital, capacity, and warranty risks.

\subsection*{72. Local Services: Geography Is Part of the Value Proposition}

Customers choosing clinics, accountants, maintenance technicians, restaurants, repair services, or local learning providers value travel time, availability, personal trust, cultural familiarity, and emergency response. A smaller service business can win through proximity and consistency rather than grand technological differentiation.

Physical access creates genuine advantage, but opening locations before measuring repeat demand creates fixed-cost exposure. The economic question is whether location adds enough customer value and revenue to repay the operational burden.

\subsection*{73. Healthcare and Consequential Credence Services}

In healthcare and similarly sensitive markets, outcomes are difficult for buyers to evaluate even after receiving the service. Trust, licensing, standards, safety, referrals, coordination, and transparency become central. Differentiation through unverified superiority claims is particularly dangerous because customers face vulnerability and severe downside.

An ethical organization differentiates through better access, comprehensible information, reliable follow-up, transparent costs, and demonstrable professional competence, remaining strictly within applicable regulation and professional duties.

\subsection*{74. Financial Services: Risk, Regulation, and Trust as the Product}

A financial-services firm may advertise better returns, faster access, or lower fees, but customers also purchase custody, risk control, regulatory compliance, financial continuity, and understandable terms. Products with identical headline yields differ substantially in risk, liquidity, fees, tax treatment, and conflicts of interest.

A new entrant differentiates through transparent total costs, better workflow integration, useful decision support, or access to underserved customers. It must avoid treating opaque risk exposure as innovation.

\subsection*{75. Deep Technology: Independent Performance Evidence}

Scientific instruments, materials, research software, robotics, and deep-tech systems derive advantage from hard-to-reproduce technical performance, integration expertise, IP, and domain know-how. Yet a laboratory result is not a purchasable product if it cannot be produced consistently, maintained, validated, or integrated into the customer's environment.

The critical proof is an application-relevant independent benchmark under representative conditions, followed by manufacturing, servicing, and cost realism. Buyers have long approval cycles and need technical reference customers.
